% !Mode:: "TeX:UTF-8"
\chapter{绪论}

\section{课题背景}

\subsection{深度学习框架的工程地位与安全隐患}

深度学习（Deep Learning, DL）框架在当前人工智能技术栈中承担底层基础设施的角色。
以PyTorch和TensorFlow为代表的主流框架，向下屏蔽异构计算硬件的物理特性与内存分配机制，
向上提供高度封装的数学算子（Operator）库与自动微分接口。
这种分层封装降低了大规模神经网络的开发门槛。
然而，框架自身已演进为由数百万行代码构成的超大型系统软件，
其内部潜伏的逻辑缺陷也在持续积累\cite{ref_dl_testing}。

与传统应用软件直接抛出异常中断不同，DL框架的缺陷往往具有隐蔽性。
在处理特定边界条件的张量维度组合或执行混合精度计算时，
底层算子可能在正向传播或梯度计算过程中产生微小的数值截断误差。
这种偏差在运行时通常不会触发系统级警告。
上层神经网络在包含逻辑错误的框架上静默完成训练后，
会将偏差以权重的形式吸收为模型自身的隐性缺陷\cite{ref_dl_safety}。
一旦此类缺陷被部署至安全关键场景——如自动驾驶的环境感知模块或医疗影像的病灶识别系统——将引发系统漏检与误判。
此外，框架底层的内存越界或缓冲区溢出漏洞也存在被攻击者利用的风险。
攻击者可通过构造包含畸形张量的计算图（Computation Graph）触发非法内存访问，
进而威胁云端算力集群的数据安全\cite{ref_dl_vulnerability}。

\subsection{模糊测试的引入与测试数据管理瓶颈}

针对上述框架深层逻辑异常，
传统基于人工编写固定输入的单元测试方法难以覆盖框架庞大的输入空间。
该空间涵盖多变的数据类型、任意的张量维度以及数百种算子的嵌套拓扑组合。
因此，学术界与工业界广泛引入了模糊测试（Fuzzing）技术\cite{ref_fuzzing_survey}。
该技术通过自动化算法高频生成包含畸形参数或极端边界状态的测试输入，
持续送入待测框架执行，以主动探索深层代码段中的潜在漏洞。

然而，高度自动化的变异引擎在连续运行期间会产出规模庞大的非结构化测试产物。
这些产物包含两部分：一是测试用例本身，
其形态为包含随机生成的高维算子组合与属性配置的代码脚本；
二是运行时产生的报错日志，
此类日志交织着应用层调用栈回溯、底层C++执行引擎的异常断点以及显卡架构的内存转储信息。
在分布式算力支持下，全天候运行的测试工具每日可产生数十万个独立脚本及数十吉字节的日志文本。
这种数据体量的膨胀与结构的无序，使得后续漏洞溯源与效能评估面临以下三方面核心痛点。

其一，缺陷回溯与根因定位存在检索困境。
系统每日捕获的数百次崩溃警告中，表层大量报错往往由底层少数几个相同漏洞引起。
传统基于关键字匹配的文本检索工具在面对包含跨语言调用、变量名随机变异的崩溃日志时，
特征匹配精度受到严重制约。
同时，底层计算引擎对指令调度的高度封装，
使得异常堆栈难以自动映射至应用层的具体触发算子，
拉长了从根因定位到缺陷修复的周期\cite{ref_fault_localization}。

其二，测试用例多样性缺乏有效的评估手段。
模糊测试的有效性依赖于生成算法对输入空间的探索广度。
变异引擎在实际运行中容易陷入局部最优状态，反复生成拓扑结构雷同的代码逻辑。
针对包含节点属性与复杂边连接关系的高维计算图数据，
传统关系型数据库与基础统计图表难以对其分布特征进行有效表征与度量\cite{ref_test_diversity}。

其三，测试效能分析的维度过于单一。
现阶段对框架模糊测试的评估常局限于"代码覆盖率"或"用例是否异常"等粗粒度标量指标。
这种模式在入库时剥离了计算图自身的深层语义特征，
且忽略了单个用例执行所消耗的时间与显存成本。
系统难以量化评估单体用例的投入产出比，也无法自动剥离低价值冗余用例\cite{ref_test_effectiveness}。
此外，现有工作流缺乏从历史失效测试集中提取算子盲区特征并将其转化为定向测试资产的闭环能力。

\section{国内外研究现状}

\subsection{深度学习框架模糊测试方法研究}

面向DL框架的模糊测试是近年来软件测试领域的研究热点。
Pei等人提出的DeepXplore\cite{ref_deepxplore}引入了神经元覆盖率的概念，
通过差分测试策略自动生成能够触发系统内部逻辑差异的测试输入，
为后续研究提供了理论基础。
Xie等人\cite{ref_deephunter}在此基础上提出了DeepHunter，
将变异策略从像素级扰动扩展至覆盖率引导的多层次变异机制。

在框架级别的测试方面，Wei等人提出的FreeFuzz\cite{ref_freefuzz}
从开源代码中自动提取API调用模式来构造测试用例，
以较低成本覆盖了PyTorch与TensorFlow的大量框架接口。
Deng等人\cite{ref_titanfuzz}进一步探索了利用大语言模型（Large Language Model, LLM）
生成框架测试用例的可行性，验证了LLM在理解框架API语义方面的潜力。
上述工作在用例自动化生成方面取得了重要进展，
但研究重心集中于变异策略与生成算法的优化，
对测试产物在生成之后的管理、归约与效能分析关注不足。
当生成端的产出规模持续膨胀时，管理端的滞后成为制约整体测试效能的瓶颈。

\subsection{测试数据管理与效能评估相关研究}

在测试用例管理领域，传统研究主要聚焦于用例优先级排序与用例约简两个方向。
Rothermel等人\cite{ref_tcp}提出了基于覆盖率的贪心排序策略，
通过优先执行覆盖新代码分支最多的用例来提升回归测试效率。
Harrold等人\cite{ref_tsm}研究了在维持覆盖率约束下的最小用例集选择问题。
这些方法在传统软件测试中已较为成熟。
但面对DL框架的测试产物时，其局限性较为明显：
传统方法处理的输入通常为标量参数或简单数据结构，
而DL测试用例的核心特征是高维计算图拓扑与多维嵌套的张量配置。
基于文本匹配或覆盖率向量的相似度计算方式，
难以捕捉此类数据的深层语义关联。

近年来，向量数据库（Vector Database）技术的发展
为高维非结构化数据的管理开辟了新的路径\cite{ref_milvus}。
以Milvus为代表的向量检索引擎支持对稠密浮点向量
执行近似最近邻（Approximate Nearest Neighbor, ANN）查询，
能够在亚秒级响应时间内完成百万级向量的相似度匹配。
将此类技术引入测试数据管理，
有望解决传统关键字检索在处理非结构化崩溃日志时精度不足的问题。
然而，将向量检索与结构化元数据存储协同应用于模糊测试数据管理的研究仍较为有限。

\subsection{大语言模型辅助软件测试研究现状}

LLM在软件工程领域的应用正受到广泛关注\cite{ref_llm_se}。
在代码生成方面，Codex\cite{ref_codex}与CodeLlama\cite{ref_codellama}等模型
展示了根据自然语言描述自动生成程序代码的能力。
在缺陷修复方面，多项研究表明LLM能够基于错误上下文生成候选修复补丁\cite{ref_llm_repair}。
这些工作为模糊测试中的用例生成与故障诊断提供了新的技术手段。

然而，通用LLM在处理垂直测试领域的任务时存在明确的局限。
脱离具体测试环境与历史运行数据，
模型在面对高度封装且语法专业的异常堆栈时，
容易产生"幻觉"（Hallucination）现象\cite{ref_llm_hallucination}。
其具体表现为虚构底层调度逻辑或推荐已废弃的参数配置。
为缓解这一问题，
检索增强生成（Retrieval-Augmented Generation, RAG）架构\cite{ref_rag}
通过在推理前从外部知识库中召回相关上下文来约束模型的输出边界，
在多个垂直领域中已取得积极效果。
但将RAG系统性地应用于DL框架的故障诊断与测试用例定向生成，
相关研究仍处于起步阶段。

\section{研究目标与主要工作}

本课题针对DL框架模糊测试产物在管理与分析环节的薄弱现状，
旨在设计并实现一套面向异构非结构化测试数据的智能管理与效能分析系统。
具体研究目标涵盖以下三个方面。

第一，构建跨模态关联映射的异构存储架构。
针对测试过程中产生的计算图脚本、异常崩溃日志及覆盖率元数据，
设计结构化元数据与高维语义向量的协同存储方案。
该方案旨在解决传统关系型数据库在处理高维语义数据时的性能瓶颈，
支撑具备语义感知能力的缺陷用例检索与归类。

第二，建立用例多样性监测与效能量化评估机制。
利用降维映射算法将高维计算图特征投射至低维平面，
实现对测试集空间探索广度的可视化监控。
同时，构建覆盖效能指数（Coverage Efficiency Index, CEI）评估模型，
量化单体用例对代码覆盖的贡献度与计算成本之比，
以此精准识别并过滤低价值冗余用例。

第三，实现基于语义增强的故障辅助诊断与闭环补盲。
整合RAG技术与LLM的推理能力，构建具备项目环境感知的辅助诊断助手。
此外，通过动态识别算子覆盖缺口驱动LLM定向生成测试脚本，
形成"特征分析—覆盖度评估—定向生成"的智能化闭环反馈流程。

本文的主要工作包括：
设计并实现自动化日志解析与双路向量化数据流水线；
部署MySQL与Milvus协同的异构混合存储底座；
开发基于深度优先搜索（Depth-First Search, DFS）拓扑回溯的故障溯源定位引擎；
提出并实现双空间Daikon故障边界定界算法；
构建基于主成分分析（Principal Component Analysis, PCA）的测试集多样性降维监测机制与CEI效能量化模型；
集成基于RAG架构的智能诊断助手与算子缺口靶向补盲模块；
在插桩编译的PyTorch环境中完成千例规模的物理对照实验，验证了系统的工程有效性。

\section{论文组织结构}

本文共分为六章，各章内容安排如下。

第一章为绪论。
本章阐述课题的研究背景与工程痛点，梳理国内外相关研究现状，明确研究目标与主要工作内容。

第二章介绍相关技术与理论基础。
本章对DL框架结构、模糊测试技术、语义编码与向量检索方法、降维分析与不变量推断理论、RAG架构以及代码覆盖率采集技术进行阐述，为后续章节提供理论支撑。

第三章进行系统需求分析与总体架构设计。
本章明确系统的功能性与非功能性需求，阐述前后端分离的混合后端架构与异构双库存储方案的设计决策，给出系统功能模块的整体划分。

第四章详细描述系统的核心算法与模块设计。
本章涵盖日志解析引擎、DFS故障溯源定位引擎、双空间Daikon边界定界引擎、CEI效能评估模型、RAG诊断助手与闭环补盲机制的算法细节。

第五章展示系统的工程实现与功能界面。
本章介绍后端服务的关键接口实现、数据库表结构设计、各功能页面的交互逻辑以及底层物理执行沙箱的构建过程。

第六章设计并执行验证实验。
本章通过模拟实验与真实物理对照实验评估CEI筛选算法的有效性，验证Daikon边界定界引擎的准确性，并对系统整体性能指标进行分析。

最后为总结与展望，回顾本课题的主要成果，分析现有系统的不足并讨论后续改进方向。
